\epanalipsi{2}

\tmima{Λεξικό του κεφαλαίου}

\renewcommand{\arraystretch}{1.5}
\begin{tabular*}{\linewidth}{@{\extracolsep{\fill}} >{\raggedright\arraybackslash}p{0.28\linewidth} >{\raggedright\arraybackslash}p{0.62\linewidth}}
\toprule
\textbf{Λέξη ή τύπος} & \textbf{Με μία πρόταση}\\
\midrule
$\dfrac{\kappa}{\nu}$ & $\kappa$ μέρη από $\nu$ ίσα μέρη, και ταυτόχρονα το $\kappa:\nu$ (με $\nu\neq 0$).\\
κλασματική μονάδα & Το $\dfrac{1}{\nu}$: ένα από τα $\nu$ ίσα μέρη.\\
ισοδύναμα & Ίδιο κομμάτι, άλλο κόψιμο: $\dfrac{\alpha}{\beta}=\dfrac{\gamma}{\delta}$ όταν $\alpha\delta=\beta\gamma$.\\
ανάγωγο & Δεν απλοποιείται άλλο: οι όροι είναι πρώτοι μεταξύ τους.\\
ομώνυμα, ετερώνυμα & Ίδιος ή διαφορετικός παρονομαστής. Για πράξεις και συγκρίσεις κάνω ετερώνυμα ομώνυμα με το ΕΚΠ.\\
μεικτός & Ολόκληρο και κλάσμα μαζί: $2\dfrac{3}{4}=2+\dfrac{3}{4}=\dfrac{11}{4}$.\\
αντίστροφος & Το κλάσμα γυρισμένο ανάποδα, με γινόμενο 1. Το 0 δεν έχει.\\
διαίρεση & Πολλαπλασιάζω με τον αντίστροφο του διαιρέτη.\\
σύνθετο κλάσμα & Κλάσμα με κλάσμα στους όρους του: η γραμμή σημαίνει διαίρεση.\\
\bottomrule
\end{tabular*}

\tmima{Ασκήσεις επανάληψης}

\ex Γράψε ως κλάσμα: α) τα πέντε από τα δώδεκα ίσα μέρη \quad β) το πηλίκο $7:9$ \quad γ) τον φυσικό 4 με παρονομαστή 6.

\ex Βρες τα $\frac{3}{8}$ του 240. Βρες και τον αριθμό που τα $\frac{5}{6}$ του είναι 75.

\ex Απλοποίησε μέχρι να γίνουν ανάγωγα: α) $\frac{36}{48}$ \quad β) $\frac{90}{150}$ \quad γ) $\frac{66}{99}$.

\ex Κάνε ομώνυμα και σύγκρινε: α) $\frac{5}{12}$ και $\frac{3}{8}$ \quad β) $\frac{7}{9}$ και $\frac{5}{6}$.

\ex Γράψε σε αύξουσα σειρά τα $\frac{1}{2}$, $\frac{5}{8}$, $\frac{3}{5}$ και $\frac{7}{10}$.

\ex Υπολόγισε: α) $\frac{3}{4}+\frac{5}{6}$ \quad β) $\frac{7}{8}-\frac{1}{3}$ \quad γ) $2\frac{1}{4}+1\frac{2}{3}$ \quad δ) $5-2\frac{3}{5}$.

\ex Υπολόγισε: α) $\frac{5}{6}\cdot\frac{9}{10}$ \quad β) $3\frac{1}{3}\cdot 2\frac{1}{4}$ \quad γ) $\frac{8}{9}:\frac{4}{3}$ \quad δ) $2\frac{1}{2}:1\frac{1}{4}$.

\ex Βρες την τιμή: α) $\left(\frac{1}{2}+\frac{1}{3}\right)\cdot\frac{6}{5}$ \quad β) $\frac{3}{4}-\frac{1}{2}:2$.

\ex Κάνε απλά τα σύνθετα κλάσματα: α) $\dfrac{\frac{3}{4}}{\frac{5}{8}}$ \quad β) $\dfrac{2+\frac{1}{2}}{1-\frac{1}{2}}$.

\ex Γράψε το $\frac{31}{6}$ ως μεικτό και τον $4\frac{3}{7}$ ως κλάσμα.

\ex \textbf{Το κέικ της Μυρτώς.} Η Μυρτώ κόβει το κέικ σε 12 ίσα κομμάτια. Η παρέα των εννέα παιδιών τρώει από ένα κομμάτι. Τι μέρος του κέικ έφαγαν και τι μέρος έμεινε; Η γιαγιά Ρήνη τρώει το μισό από αυτό που έμεινε. Τι μέρος ολόκληρου του κέικ είναι αυτό;

\ex \textbf{Οι δραστηριότητες.} Η Α3 έχει 27 μαθητές. Τα $\frac{4}{9}$ παίζουν μπάσκετ και το $\frac{1}{3}$ μαθαίνουν μουσική (κανείς δεν κάνει και τα δύο). Πόσα παιδιά κάνουν το καθένα και τι μέρος της τάξης δεν κάνει καμία από τις δύο;

\ex \textbf{Η προπόνηση.} Ο Νίκος έτρεξε το πρωί τα $\frac{2}{5}$ μιας διαδρομής και το απόγευμα τα $\frac{3}{10}$. Τι μέρος έτρεξε συνολικά και τι μέρος του μένει; Αν η διαδρομή είναι 20 χλμ., πόσα έτρεξε;

\ex \textbf{Το διάβασμα.} Ο Βασίλης χρειάζεται $\frac{3}{4}$ της ώρας για κάθε μάθημα. Πόσα μαθήματα προλαβαίνει σε 3 ώρες;

\ex \textbf{Ερωτήσεις αυτοαξιολόγησης (έννοια, ισοδύναμα).} Σωστό ή λάθος; α) Το $\frac{5}{5}$ είναι ίσο με 1. \quad β) Κάθε κλάσμα έχει αριθμητή μικρότερο από τον παρονομαστή. \quad γ) $\frac{4}{6}=\frac{2}{3}$. \quad δ) Το $\frac{3}{7}$ είναι ανάγωγο. \quad ε) Αν προσθέσω 2 και στον αριθμητή και στον παρονομαστή, το κλάσμα μένει ίδιο. \quad στ) $\frac{0}{9}=0$.

\ex \textbf{Ερωτήσεις αυτοαξιολόγησης (σύγκριση, πράξεις).} Σωστό ή λάθος; α) $\frac{3}{8}>\frac{3}{5}$. \quad β) $\frac{1}{2}+\frac{1}{2}=\frac{2}{4}$. \quad γ) $\frac{2}{3}\cdot\frac{3}{2}=1$. \quad δ) $6:\frac{1}{3}=2$. \quad ε) $1\frac{3}{4}=\frac{7}{4}$. \quad στ) Το άθροισμα δύο κλασμάτων είναι μεγαλύτερο από καθένα τους. \quad ζ) Το γινόμενο δύο κλασμάτων μικρότερων από το 1 είναι μικρότερο από καθένα τους.

\ex \textbf{Ερωτήσεις αυτοαξιολόγησης (αντίστροφοι, διαίρεση).} Σωστό ή λάθος; α) Ο αντίστροφος του $\frac{2}{5}$ είναι ο $\frac{5}{2}$. \quad β) Ο αντίστροφος του 0 είναι το 0. \quad γ) $\frac{1}{2}:\frac{1}{4}=2$. \quad δ) Η διαίρεση με κλάσμα μικρότερο από το 1 μικραίνει τον αριθμό. \quad ε) $3:\frac{3}{4}=4$. \quad στ) $\frac{5}{6}:\frac{5}{6}=1$.

\tmima{Ο χάρτης της επανάληψης}

Μετά τις λύσεις, δες σε ποιες ασκήσεις δυσκολεύτηκες. Ο πίνακας σου λέει ποια παράγραφο να ξαναδιαβάσεις.

\renewcommand{\arraystretch}{1.5}
\begin{tabular*}{\linewidth}{@{\extracolsep{\fill}} >{\raggedright\arraybackslash}p{0.38\linewidth} >{\raggedright\arraybackslash}p{0.52\linewidth}}
\toprule
\textbf{Αν δυσκολεύτηκες στις} & \textbf{ξαναδιάβασε τις παραγράφους}\\
\midrule
Ε2.1, Ε2.2, Ε2.12 & 2.1 (έννοια του κλάσματος)\\
Ε2.3 & 2.2 (ισοδύναμα κλάσματα)\\
Ε2.4, Ε2.5 & 2.2 και 2.3 (ομώνυμα, σύγκριση)\\
Ε2.6, Ε2.10, Ε2.13 & 2.4 (πρόσθεση, αφαίρεση, μεικτοί)\\
Ε2.7 (α, β), Ε2.8, Ε2.11 & 2.5 (πολλαπλασιασμός)\\
Ε2.7 (γ, δ), Ε2.9, Ε2.14 & 2.6 (διαίρεση)\\
Ε2.15, Ε2.16, Ε2.17 & 2.1 έως 2.6\\
\bottomrule
\end{tabular*}
